% mage-auto-heading: semantic-commitment
% mage-auto-heading: representing-and-interacting-with-models
% mage-auto-heading: from-analysis-to-engineering-decisions
% mage-auto-heading: the-engineering-modeling-repertoire
% mage-auto-heading: semantic-commitment-and-analysis
% Authoritative LaTeX manuscript; Markdown is a generated compatibility projection.

% <!-- part-title: Modeling -->

% <!-- chapter-title: Constructing and Using Engineering Models -->

\chapter{Constructing and Using Engineering Models}

% <!-- section-terms: modeling-principle, degrees-of-freedom -->

% <!-- point: the-direction-also-runs-in-reverse | An exploratory implementation can expose constraints, tradeoffs, state distinctions, dependencies, and bounds that were not evident in the initial design. | terms: modeling-principle, degrees-of-freedom -->
% <!-- point: human-authority-is-not-human-invention | The agent may supply a design insight; the engineer takes responsibility for adopting it by evaluating it against the system's requirements. | terms: modeling-principle -->

We can model a system before we build it, after we build it, or while we build it. A requirement can supply an obligation before any code exists. An existing system and the people who maintain it can supply observed dependencies and the reasons for preserving some of them. An agent building a prototype may discover a state distinction that nobody thought to specify. The engineer can start from any of the sources in \autoref{multiple-routes-to-engineering-models} and judge what each establishes about the system, including which properties later implementations must preserve.

% <!-- point: model-origins-and-uses | Models support engineering design, direct agent work, and supply properties against which results can be checked. | terms: modeling-principle -->

However a model originates, its value depends on what engineers can do with it. Consider the three uses beneath the model in \autoref{multiple-routes-to-engineering-models}. An engineer can reason through a proposed service boundary by asking which clients would need to change if the database schema changed. The engineer can direct an agent adding a client to use the interface named in the model. A reviewer can then compare the client's actual dependencies with the permitted ones. Detecting an unexpected dependency gives the reviewer evidence; rejecting a change that violates the boundary requires an additional enforcement mechanism.

\begin{figure}[htbp]
\centering
\includegraphics[width=\linewidth]{book/latex/assets/multiple-routes-to-engineering-models-abe57356fe.pdf}
\caption{\emph{Multiple routes to engineering models.} A model may be constructed from requirements and intent, recovered from existing systems and experience, or discovered through agent exploration. However it arrives, the model carries engineering knowledge and the selected obligations drawn from it, and it can support three uses: engineers reason and design through it, agents are directed by it, and results are evaluated and checked against it. A model recovered from an implementation describes what exists; engineers judge whether its design is desirable and which discovered properties should become obligations. Further engineering discoveries may motivate new or revised models.}
\label{multiple-routes-to-engineering-models}
\end{figure}

% <!-- point: the-most-important-evidence-comes-from-realization | Implementing and trying a design can reveal distinctions that were not evident in the initial model. | terms: modeling-principle -->
% <!-- point: revise-the-model-when-evidence-changes-the-judgment | Neither direction guarantees the boundary is right; analysis, implementation, measurement, or failure may reveal that an omitted distinction mattered or an included one did not, and the model is revised when the evidence changes the engineering judgment. | terms: degrees-of-freedom, model-as-map -->

Models also change as engineers learn more about a system. Measurements may reveal that a performance model omits an important source of latency; a prototype may expose a state or interaction that the original requirements did not anticipate. Engineers and agents can use this evidence to revise the model, adding distinctions needed for new questions or removing details that no longer serve its purpose. Modeling continues throughout development as new evidence changes what engineers need to understand and preserve.

\section{The Modeling Repertoire for Engineering}\label{heading:the-engineering-modeling-repertoire}

% <!-- section-terms: model-classes, projection, modeling-principle -->

% <!-- point: engineering-fields-have-accumulated-a-modeling-repertoire | Software engineers need not invent a representation from first principles each time; engineering fields have accumulated a large repertoire of models for recurring questions, each preserving different information and supporting different analyses. | terms: model-classes, modeling-principle -->
% <!-- point: recognize-the-question-select-the-representation | The practical skill is recognizing the engineering question and selecting a representation that preserves the properties needed to answer it; MAGE prescribes no new notation. | terms: modeling-principle, model-classes -->
% <!-- point: reuse-the-established-repertoire-choosing-is-judgment | This chapter illustrates modeling practice rather than exhausting it; once the engineering question is known, start with the representations the field has already developed — the repertoire is reusable engineering knowledge, and choosing from it remains engineering judgment. | terms: modeling-principle, model-classes -->
% <!-- index-def: model-classes -->

The engineer's role is to identify the engineering question and select a model that makes it easier to answer. Engineers do not need to invent a modeling technique from first principles for each question, and generally should not. Decades of work in software engineering, systems engineering, mathematics, and operations research have produced models of structure, behavior, decisions, quantities, and evidence. MAGE leaves the choice of notation to the engineer. \autoref{modeling-repertoire} illustrates the available repertoire with nine examples, showing the properties each model can express and the analyses it supports.

We will organize our study of these models around three kinds of engineering question: how a system is organized, how it behaves, and how we should choose among alternatives. These are not mutually exclusive classes of representation. A quantitative model, for example, might describe a system’s resource requirements, predict its behavior, or help us choose a design. We will use such models alongside the structural, behavioral, and decision models developed in the following sections.

% <!-- point: read-the-repertoire-as-available-moves-not-a-taxonomy | What matters is not a representation's graphical shape but the engineering question its entities and relations preserve. | terms: model-classes, projection -->
% <!-- point: shape-does-not-distinguish-a-model-semantics-do | Shape does not distinguish representations, semantics do; semantic commitment decides how much of that meaning is declared rather than understood. | terms: system-knowledge-graph, semantic-commitment -->

\begin{figure}[htbp]
\centering
\includegraphics[width=\linewidth]{book/latex/assets/modeling-repertoire-f010848f36.pdf}
\caption{\emph{The engineering modeling repertoire.} Familiar representations preserve different relationships for different engineering questions. Similar graphical forms can carry different semantics: an arrow may represent a dependency, transition, flow, permission, ownership relation, or derivation.}
\label{modeling-repertoire}
\end{figure}

How do you choose among these representations? Start with what you need to establish about the system. A requirement that only one worker may own a job raises a question about concurrent claims. To answer it, explore whether any permitted sequence of claims and releases leaves two workers holding the same job. The model must therefore represent ownership and the transitions that change it. Failures and proposed changes can raise questions too. \autoref{engineering-questions-and-candidate-models} pairs each analysis with models that supply the facts it needs.

For a simple implementation, we may be able to answer the question by inspecting the code. A separate model becomes attractive when the relevant facts are scattered across the implementation, when the possible interactions are difficult to examine directly, or when we expect to ask related questions repeatedly.

{\footnotesize
\begin{longtable}[]{@{}
  >{\raggedright\arraybackslash}p{(\linewidth - 4\tabcolsep) * \real{0.18}}
  >{\raggedright\arraybackslash}p{(\linewidth - 4\tabcolsep) * \real{0.54}}
  >{\raggedright\arraybackslash}p{(\linewidth - 4\tabcolsep) * \real{0.28}}@{}}
\caption[Engineering questions, analyses, and candidate models]{\emph{Engineering questions, analyses, and candidate models.} The intended analysis determines what a useful model must represent.}\label{engineering-questions-and-candidate-models}\tabularnewline
\toprule\noalign{}
Engineering concern & Question and analysis & Candidate models \\
\midrule\noalign{}
\endhead
\bottomrule\noalign{}
\endlastfoot
Correctness & Can two workers own the same job? Explore concurrent states and check the ownership invariant. & Ownership and state-transition models \\
Performance & Will requests accumulate faster than they can be served? Compare arrival rates with service capacity. & Queueing and workload models \\
Security & Can an untrusted component reach protected data through an unauthorized path? Trace dependencies and check permissions along each path. & Dependency graph and trust-boundary model with permissions \\
Reliability & Can a failed operation recover without losing work? Explore failure and recovery sequences. & Recovery state machine \\
Regulatory compliance & Which stores and processors are implicated by a GDPR erasure request? Trace where personal data was sent and which records were derived from it. & Data-flow and provenance models \\
\end{longtable}
}

We can use the same model to answer different questions. The same state machine can help answer questions about correctness and reliability. We can examine its transitions for conflicting claims, then follow its recovery paths to see whether work survives a failure. Some questions require combining models because no single representation preserves all the relevant facts.

% <!-- box-family: inset -->

% <!-- inset-domain: ENGINEERING -->

% mage-box-prefix: Inset —

% <!-- print-keep-with-next: 6 -->

\begin{bookbox}{inset}{Debugging Models Before Code}
\label{inset-debugging-models-before-code}
One engineering question may require several models. Suppose we are designing a battery-powered device with two sensors that collect measurements and transmit them over a radio. Can any measurements be lost? \autoref{sensor-buffer-radio} answers this question by combining three models:

\begin{itemize}
\tightlist
\item \textbf{(a) Structural model.} Two sensors feed a shared buffer in on-device storage, from which a radio transmits measurements. The model establishes that both sensors compete for the same storage capacity.
\item \textbf{(b) Quantitative model.} Each sensor produces one measurement per second, the buffer holds three measurements, and the radio transmits every five seconds. These parameters do not specify how many sensors share the buffer.
\item \textbf{(c) Behavioral model.} An arriving measurement is stored if space remains and discarded otherwise. A transmission sends and removes all stored measurements; when transmission and sampling coincide, transmission occurs first. Starting with an empty buffer, two measurements arrive at second 1. At second 2, only one of the next two fits; the other is lost. The derived execution trace combines the structure, quantities, and behavioral rules to expose the failure.
\end{itemize}

We can now consider increasing buffer capacity, transmitting more frequently, or sampling less often. Alternatively, we could change the overflow behavior to retain measurements elsewhere until transmission, provided that additional storage and eventual transmission capacity are sufficient. Each choice affects energy consumption, memory requirements, measurement quality, or implementation complexity. I would much rather discover this problem by examining three small models than by debugging a thousand lines of C.

\begin{figure}[htbp]
\centering
\includegraphics[width=\linewidth,alt={Two sensors share a three-measurement buffer in on-device storage. Each samples once per second; the radio transmits every five seconds. Arrivals are stored if space remains and discarded otherwise. The derived trace shows one measurement lost at second two. At second five, transmission empties the buffer before two new measurements arrive.}]{book/latex/assets/sensor-buffer-radio-de7cbf9651.pdf}
\caption[Combining three models to expose measurement loss]{\emph{Combining three models to expose measurement loss.} The structural model identifies two sensors sharing on-device storage; the quantitative model specifies per-sensor sampling rates, buffer capacity, and transmission interval; and the behavioral model specifies what happens when measurements arrive at a full buffer. Their combination reveals that a measurement is lost at second 2, before the first transmission.}
\label{sensor-buffer-radio}
\end{figure}
\end{bookbox}

\section{Semantic Commitment and Analysis}\label{heading:semantic-commitment-and-analysis}

% <!-- section-terms: modeling-principle, semantic-commitment, system-knowledge-graph, structured -->

\subsection{Semantic Commitment}\label{heading:semantic-commitment}

% <!-- point: a-model-can-preserve-distinctions-without-defining-them | A representation can preserve the right distinctions while its interpreters supply most of its meaning. | terms: semantic-commitment, model-as-map -->
% <!-- point: agents-make-the-semantics-tradeoff-newly-consequential | Explicit semantics preserve meanings that later agents would otherwise have to infer without the shared context of the original team. | terms: semantic-commitment, externalized-knowledge -->

Engineering models range from whiteboard sketches to specifications written in a programming language. A team sketching a service boundary can agree in conversation that an arrow means one service is permitted to call another. That agreement may be precise enough for the people at the whiteboard. An agent reading the sketch later, or a tool checking a proposed dependency, needs some way to recover the same interpretation. How much meaning the engineer needs to record depends on who will use the model and what they need to do with it.

% <!-- point: semantic-commitment-names-the-second-axis | Semantic commitment: how much meaning the representation itself carries explicitly. | terms: semantic-commitment, structured -->
% <!-- point: semantic-commitment-is-distinct-from-purposeful-reduction | Purposeful reduction determines which distinctions to preserve; semantic commitment determines how precisely the representation defines them. | terms: semantic-commitment, model-as-map -->
% <!-- index-def: semantic-commitment -->

\textbf{Semantic commitment} is how much of a representation's meaning it defines explicitly and how much its interpreter must supply. A diagram file can preserve every arrow while leaving its meaning ambiguous. A small typed graph, on the other hand, can define exactly which calls are permitted while leaving every service's algorithm open. Defining relations, units, and constraints can make stronger analyses possible, but the definitions take work to write, learn, and maintain. \autoref{semantic-commitment-spectrum} compares what several representations define and what each leaves unspecified.

The choice also has consequences for Alignment. A representation with sufficiently defined semantics may allow us to derive implementation artifacts and checks mechanically, rather than ask an engineer or agent to interpret its meaning each time. \autoref{3.2-governance-mechanisms:heading:semantic-commitment-and-mechanization} examines this opportunity.

% <!-- point: semantic-precision-is-not-generative-completeness | Semantic precision and generative completeness are different properties. | terms: semantic-commitment, degrees-of-freedom -->
% <!-- point: stronger-semantics-buy-capabilities-and-create-obligations | Stronger semantics buy capabilities and create obligations; they are not inherently better. | terms: semantic-commitment, degrees-of-freedom -->
% <!-- point: the-spectrum-in-concrete-artifacts | Concrete artifacts illustrate the range; the rows are examples, not maturity levels. | terms: semantic-commitment, degrees-of-freedom -->
% <!-- point: the-columns-are-not-one-progression | Stronger commitment aids mechanical interpretation; it moves interpretation burden rather than removing it. | terms: semantic-commitment -->

% <!-- print-keep-with-next: 14 -->

{\footnotesize
\begin{longtable}[]{@{}
  >{\raggedright\arraybackslash}p{(\linewidth - 4\tabcolsep) * \real{0.23}}
  >{\raggedright\arraybackslash}p{(\linewidth - 4\tabcolsep) * \real{0.39}}
  >{\raggedright\arraybackslash}p{(\linewidth - 4\tabcolsep) * \real{0.38}}@{}}
\caption[Degrees of semantic commitment, in concrete artifacts]{\emph{Semantic commitment is not model completeness.} Purposeful reduction determines which distinctions a model preserves. Semantic commitment determines how precisely it defines them. UML 2.5.1~\citep{omg2017uml} and SysML v2~\citep{omg2025sysml} with KerML 1.0~\citep{omg2025kerml} define standard meanings for selected constructs.}\label{semantic-commitment-spectrum}\tabularnewline
\toprule\noalign{}
Representation & Meaning available to the reader or tool & What remains unspecified \\
\midrule\noalign{}
\endhead
\bottomrule\noalign{}
\endlastfoot
Informal architecture sketch & Labels identify the components. & An unlabeled arrow may denote an observed call or permission to call. \\
UML 2.5.1 state machine & States, triggers, guards, and transitions have defined behavioral meanings. & A lifecycle can leave the implementation of its event handlers open. \\
OpenAPI 3.1, JSON Schema, or typed repository registries & Named fields, types, constraints, and references support queries and validation. & A declared field type alone leaves its business meaning and storage unspecified. \\
SysML v2 with KerML 1.0 & Types, relationships, expressions, and constraints support composition and analysis. & A system model can fix a component boundary while leaving the internal design open. \\
\end{longtable}
}

% <!-- point: state-machine-semantics-make-path-properties-checkable | An explicit transition relation allows a path property to be checked; composing machines also requires a choice of communication semantics. | terms: semantic-commitment, model-as-map -->

% <!-- box-family: inset -->

% <!-- inset-domain: FORMAL METHODS -->

% mage-box-prefix: Inset —

\begin{bookbox}{inset}{What Does a State-Machine Arrow Mean?}
\label{inset-state-machine-arrow}
Find the panel labeled “Finite-state” in \autoref{modeling-repertoire}. It shows four job states: IDLE, RUNNING, DONE, and FAILED. Suppose a job starts in IDLE, each arrow denotes a permitted transition, and DONE and FAILED are terminal. We can now ask whether a job can reach DONE without first entering RUNNING.

Let \(\operatorname{Paths}(M)\) be the paths in model \(M\) that start at IDLE and follow its arrows, and let \(p_i\) denote the state at position \(i\) on a path \(p\), counting from zero. For every position on every path, the requirement is:
\[
\forall p \in \operatorname{Paths}(M),\quad
p_i = \mathrm{DONE} \;\Longrightarrow\;
\exists j < i:\ p_j = \mathrm{RUNNING}.
\]
A search can check the property by tracking whether RUNNING has been visited. In the diagram, the only arrow into DONE comes from RUNNING. A job can still end in FAILED.

Now compose the job model with a worker that reports completion. Under synchronous semantics, the worker reports completion and the job enters DONE in one joint transition. Under asynchronous semantics, the worker sends a message and the job processes it later. Ownership might change in between. The job's arrows suffice to check the local path property. To check whether a former owner's completion can be accepted, we must also specify how the two machines communicate.
\end{bookbox}

\subsection{Representing and Interacting with Models}\label{heading:representing-and-interacting-with-models}

% <!-- point: structured-is-the-word-doing-the-work | Structure — an explicit declared shape a machine can read — is what lets a representation be traversed, checked, and joined. | terms: structured, system-knowledge-graph -->
% <!-- point: the-simplest-useful-form-is-a-graph | The simplest useful form of externalized knowledge is a graph of engineering entities and their relations; graph names the structure, not the storage. | terms: system-knowledge-graph, structured -->
% <!-- index-def: structured -->
% <!-- index-def: system-knowledge-graph -->

Models can be maintained in ordinary project files, databases, or specialized modeling environments. The choice resembles deciding whether to store application data in flat files, use SQLite, or operate a database server. Specialized modeling tools may provide graphical editing, consistency checks, simulation, and traceability across models. Interfaces using the Model Context Protocol (MCP) can expose these modeling capabilities as tools an agent can invoke. The appropriate choice depends on the analyses required, the team's familiarity with the tools, and the cost of learning, integrating, and maintaining them.

% <!-- point: capability-comes-in-rungs | Increasing semantic commitment can make additional capabilities possible — not a required sequence or maturity model. | terms: capability-ladder, structured, semantic-commitment -->
% <!-- index-def: capability-ladder -->

Whatever the tooling, models become more useful when engineers and agents can query their elements and relationships. A service registry, for example, can identify each service, its owner, dependencies, and deployment. An agent investigating a service can retrieve the relevant records without reconstructing those relationships from source code, even when the complete model exceeds its working context. For example, Spotify makes service ownership and documentation available to agents through its service catalog~\citep{spotify2026constraint}. Other models support operations beyond lookup: a schema can validate a document or generate interface types, while a transition model can drive a simulation or direct a running system. These operations depend on the meaning assigned to the model's elements, not merely on their availability.

\subsection{From Analysis to Engineering Decisions}\label{heading:from-analysis-to-engineering-decisions}

% <!-- point: validate-the-model-for-its-question | Validation examines whether the model represents the distinctions, assumptions, and behaviors needed for the question. | terms: semantic-commitment, model-as-map -->

% <!-- point: analysis-produces-evidence-for-an-engineering-decision | A model's semantics determine which analyses can produce evidence for a design decision. Implementation conformance requires separate evidence. | terms: semantic-commitment, modeling-principle -->

\autoref{question-to-consequence-progression} follows an engineering question from model selection through analysis to a decision. Suppose an engineer wants to know whether a recovery policy can lose a job. Finding a path to loss in the model gives the engineer a sequence to investigate and may warrant changing the policy. If the search finds no such path but a trace shows that a job was lost, the trace can expose a missing failure transition. The engineer can revise the model and repeat the search. Adopting the resulting recovery rule is a design decision. Establishing that the implementation follows it takes evidence from the implementation.

\begin{figure}[htbp]
\centering
\includegraphics[width=\linewidth]{book/latex/assets/question-to-consequence-progression-25598a02ec.pdf}
\caption{\emph{From engineering question to engineering consequence.} The representation determines which analyses become tractable. Their results can inform a design decision, expose a property for checking, or both. Chapter 3 takes up the separate question of which obligations should be enforced.}
\label{question-to-consequence-progression}
\end{figure}

The following sections examine these choices for three familiar classes of engineering question. Structural models preserve the relationships needed to ask what is connected to what. Behavioral models preserve states and transitions so we can ask what may happen. Decision models preserve the criteria for choosing what should happen next. A single system can raise all three questions, and their answers can draw on the same entities.
